\chapter{Método} \label{chp:method}

%%%%%%%%%%%%%%%%%%%%%%%%%%%%%%%%%%%%%%%%%%%%%%%%%%%%%%%%%%%%%%%%%%%%%%%%%%%%%%%%
%%%%%%%%%%%%%%%%%%%%%%%%%%%%%%%%%%%%%%%%%%%%%%%%%%%%%%%%%%%%%%%%%%%%%%%%%%%%%%%%

\section{Preparación de los datos}

Se separó \texttt{row_id} de los predictores y de la variable objetivo. Las
variables categóricas se transformaron con \texttt{OneHotEncoder} dentro de un
flujo de modelado; las numéricas se mantuvieron sin escalado. El conjunto de
test no se utilizó durante el desarrollo.

La evaluación empleó 5-fold stratified cross-validation, manteniendo
aproximadamente la misma proporción de clases en cada fold. Se
registraron accuracy, balanced accuracy, F1 y ROC-AUC. La ROC-AUC se utilizó
para comparar configuraciones sin depender de un umbral concreto.

\section{Algoritmo de clasificación}

Se utilizó un árbol de decisión CART para clasificación, implementado mediante
\texttt{DecisionTreeClassifier} de \texttt{scikit-learn}. CART construye el
modelo mediante particiones binarias sucesivas, en cada nodo selecciona una
variable y un umbral que separan las observaciones en dos ramas. El criterio
utilizado fue Gini, que es la configuración por defecto del clasificador y
mide la impureza de cada partición. Este algoritmo resulta adecuado para la
Phase I porque sus decisiones pueden expresarse directamente como reglas y no
requiere escalar las variables numéricas.

Se optó por CART y no por un Generalized Optimal Sparse Decision Tree (GOSDT).
Aunque GOSDT busca árboles óptimos y dispersos, está pensado para variables
discretas, mientras que este problema incluye predictores numéricos.

\section{Selección de la complejidad}

Como referencia se ajustó primero un árbol sin restricciones.

A continuación se exploraron distintas profundidades máximas y tamaños mínimos
de hoja, con el objetivo de identificar un compromiso razonable entre rendimiento
predictivo e interpretabilidad. En concreto, se recorrió un rango de
profundidades y varios tamaños mínimos de hoja, comparando las métricas de
validación cruzada con el tamaño resultante del árbol. La selección prioriza
estructuras simples cuando su rendimiento es comparable al de alternativas más
complejas.

\section{Selección de variables}

Con la complejidad fijada se realizó una selección forward de predictores,
partiendo de un conjunto vacío. En cada paso se añadió la variable que más
mejoraba la ROC-AUC media y se escogió el menor subconjunto cuyo resultado
quedaba dentro de un error estándar del mejor. Cada candidato se evaluó con el
mismo 5-fold stratified cross-validation, ajustando el preprocesado dentro de
cada fold para evitar usar datos de validación. Cabe destacar que la estimación
de esta selección es ligeramente optimista porque selección y evaluación
comparten los folds.

\section{Umbral y explicaciones locales}

Después de fijar la configuración, el árbol se reentrenó con el conjunto de
entrenamiento. El umbral se eligió a partir de predicciones out-of-fold,
buscando maximizar la balanced accuracy. Para las explicaciones locales se
seleccionó un caso correctamente clasificado de cada clase.
